\documentclass[tikz,border=8pt]{standalone}
\usepackage{fontspec}
\usetikzlibrary{arrows.meta}

\setmainfont{Arial}
\setsansfont{Arial}

\definecolor{pagebg}{HTML}{F8FAFC}
\definecolor{border}{HTML}{CBD5E1}
\definecolor{textmain}{HTML}{0F172A}
\definecolor{textmuted}{HTML}{475569}
\definecolor{blue}{HTML}{2563EB}
\definecolor{bluefill}{HTML}{DBEAFE}
\definecolor{violet}{HTML}{7C3AED}
\definecolor{cyan}{HTML}{0891B2}
\definecolor{green}{HTML}{16A34A}
\definecolor{orange}{HTML}{EA580C}
\definecolor{red}{HTML}{DC2626}
\definecolor{darkband}{HTML}{111827}

\begin{document}
\begin{tikzpicture}[x=1mm,y=1mm,line cap=round,line join=round,font=\sffamily]
  \fill[pagebg] (0,0) rectangle (360,240);
  \draw[fill=white,draw=border,line width=0.55mm,rounded corners=4mm]
    (6,6) rectangle (354,234);

  \node[font=\sffamily\bfseries\fontsize{18}{22}\selectfont,text=textmain]
    at (180,217) {Линейчатые спектры и уровни энергии};
  \node[font=\sffamily\fontsize{10.5}{13}\selectfont,text=textmuted]
    at (180,201) {Дискретные энергии атома приводят к отдельным линиям спектра};

  % Спектр уровней энергии атома водорода
  \draw[fill=pagebg,draw=border,line width=0.45mm,rounded corners=2.5mm]
    (14,45) rectangle (226,185);
  \node[font=\sffamily\bfseries\fontsize{12.5}{15}\selectfont,text=blue]
    at (120,173) {Спектр уровней энергии атома водорода};

  \draw[-{Stealth[length=3.4mm,width=2.4mm]},draw=textmuted,line width=0.65mm]
    (40,58) -- (40,160);
  \node[font=\sffamily\bfseries\fontsize{10}{12}\selectfont,text=textmuted]
    at (40,166) {$E$};

  \draw[draw=textmain,line width=0.9mm] (70,61) -- (178,61);
  \node[anchor=east,font=\sffamily\bfseries\fontsize{9}{11}\selectfont,text=textmain]
    at (65,61) {$n=1$};
  \node[anchor=west,font=\sffamily\fontsize{8.8}{10.5}\selectfont,text=textmuted]
    at (183,61) {$-13{,}6$ эВ};

  \draw[draw=textmain,line width=0.9mm] (70,111) -- (178,111);
  \node[anchor=east,font=\sffamily\bfseries\fontsize{9}{11}\selectfont,text=textmain]
    at (65,111) {$n=2$};
  \node[anchor=west,font=\sffamily\fontsize{8.8}{10.5}\selectfont,text=textmuted]
    at (183,111) {$-3{,}4$ эВ};

  \draw[draw=textmain,line width=0.9mm] (70,135) -- (178,135);
  \node[anchor=east,font=\sffamily\bfseries\fontsize{9}{11}\selectfont,text=textmain]
    at (65,135) {$n=3$};
  \node[anchor=west,font=\sffamily\fontsize{8.8}{10.5}\selectfont,text=textmuted]
    at (183,135) {$-1{,}51$ эВ};

  \draw[draw=textmain,line width=0.9mm] (70,149) -- (178,149);
  \node[anchor=east,font=\sffamily\bfseries\fontsize{9}{11}\selectfont,text=textmain]
    at (65,149) {$n=4$};
  \node[anchor=west,font=\sffamily\fontsize{8.8}{10.5}\selectfont,text=textmuted]
    at (183,149) {$-0{,}85$ эВ};

  \draw[draw=textmuted,line width=0.65mm,densely dashed] (70,163) -- (178,163);
  \node[anchor=east,font=\sffamily\bfseries\fontsize{9}{11}\selectfont,text=textmain]
    at (65,163) {$n\to\infty$};
  \node[anchor=west,font=\sffamily\fontsize{8.8}{10.5}\selectfont,text=textmuted]
    at (183,163) {$0$ эВ};

  % Несколько разрешённых переходов
  \draw[-{Stealth[length=3.5mm,width=2.5mm]},draw=violet,line width=1.0mm]
    (104,145) -- (104,115);
  \draw[-{Stealth[length=3.5mm,width=2.5mm]},draw=cyan,line width=1.0mm]
    (132,131) -- (132,115);
  \draw[-{Stealth[length=3.5mm,width=2.5mm]},draw=red,line width=1.0mm]
    (158,107) -- (158,65);

  \node[font=\sffamily\fontsize{8.5}{10.5}\selectfont,text=textmuted]
    at (120,50) {вертикальные расстояния показаны условно};

  % Оптический линейчатый спектр
  \draw[fill=white,draw=border,line width=0.45mm,rounded corners=2.5mm]
    (236,45) rectangle (346,185);
  \node[font=\sffamily\bfseries\fontsize{12.5}{15}\selectfont,text=blue]
    at (291,169) {Линейчатый спектр света};
  \node[font=\sffamily\fontsize{9}{11}\selectfont,text=textmuted]
    at (291,154) {отдельные линии вместо сплошной полосы};

  \draw[fill=darkband,draw=darkband,rounded corners=1.5mm]
    (247,108) rectangle (335,143);
  \draw[draw=violet,line width=1.25mm] (258,111) -- (258,140);
  \draw[draw=blue,line width=1.25mm] (266,111) -- (266,140);
  \draw[draw=cyan,line width=1.25mm] (279,111) -- (279,140);
  \draw[draw=green,line width=1.25mm] (297,111) -- (297,140);
  \draw[draw=orange,line width=1.25mm] (317,111) -- (317,140);
  \draw[draw=red,line width=1.25mm] (328,111) -- (328,140);

  \node[align=center,font=\sffamily\bfseries\fontsize{9.2}{11.5}\selectfont,text=textmain]
    at (291,86) {Один разрешённый переход\\между уровнями — одна линия};
  \node[align=center,font=\sffamily\fontsize{8.2}{10}\selectfont,text=textmuted]
    at (291,60) {Цвета и расстояния между\\линиями показаны условно};

  \draw[fill=bluefill,draw=blue!45,line width=0.45mm,rounded corners=2mm]
    (14,14) rectangle (346,36);
  \node[font=\sffamily\bfseries\fontsize{9.7}{12}\selectfont,text=blue]
    at (180,25) {Уровни описывают энергию атома, линии — энергию излучённых или поглощённых фотонов};
\end{tikzpicture}
\end{document}
